\subsection[The description of $\hat\Phi^{*k}$]{The description of $\boldsymbol{\hat\Phi^{*k}}$}\label{ssec:convolpower}

Now we are ready to describe $\tilde\rho(t\otimes a_1\cdots a_n)$ appearing in~\eqref{eq:rhoiota}, where $t$ is a Schr\"oder tree with $n+1$ leaves and elements $a_1,\dots, a_n\in \A$. First, we will need some notation.

Let $t$ be a Schr\"oder tree with $n+1$ leaves, $w=a_1\cdots a_n\in T_+(\A)$ and $k \in \{1,\dots, n\}$. Our objective is to compute $\hat\Phi^{*k}(t\otimes a_1\cdots a_n)$. By definition of the convolution product, we have that\looseness=-1
\begin{equation}
\label{eq:iterated}
	\hat\Phi^{*k}(t\otimes w) = m_{\mathbb{C}}^{[k]} \circ \hat\Phi^{\otimes k}\circ \delta_\A^{[k]}(t\otimes w),
\end{equation}
where $\delta_\A^{[k]}$ is the $k$-fold iterated coproduct
\[
	\delta_\A^{[k]} = \big(\delta_\A \otimes \Id_{\cHS(\A)}^{\otimes k-2}\big)\circ \delta_{\A}^{[k-1]},
\]
with $\delta_\A^{[2]} = \delta_\A$ and $\delta_\A^{[1]} = \Id_{\cHS(\A)}$. A similar notation applies to $m_{\mathbb{C}}^{[k]}$.

 We will also need the following combinatorial object associated to a Schr\"oder tree.

\begin{Definition}
Let $t$ be a Schr\"oder tree with $m$ internal vertices. We define the \textit{skeleton of $t$}, denoted by $\sk(t)$, to be the planar rooted tree with $m$ vertices each of them in correspondence with a unique internal vertex of $t$, and such that there is an arrow from $v$ to $w$ in $\sk(t)$ if and only if, the corresponding vertex to $w$ in $t$ is a child of the corresponding vertex to $v$ in $t$.
\end{Definition}

\begin{Example}
Let $t$ be the unlabelled Schr\"oder tree pictured in Figure \ref{fig:DecSchTree}. Its skeleton $\sk(t)$ is given by the tree
\begin{tikzpicture}[scale=0.4,
level 1/.style={level distance=7mm,sibling distance=12mm}]
\node(0)[solid node,label=above:{}]{}
child{node(1)[solid node]{}
child{[black] node(11)[solid node]{}}
}
child{node(2)[solid node]{}
child{[black] node(21)[solid node]{}}
};
\end{tikzpicture}.
\end{Example}

\begin{Remark}\label{rem:generalcuts}For any planar rooted tree $s$, Definition~\ref{def:admcut} can be extended to the notion of \textit{general admissible cut on $s$} as a subset $c$ of the set of vertices of $s$ such that any path from the root to any leaf, there is at most one vertex of the path contained in $c$. Observe that the difference with Definition~\ref{def:admcut} is that we are allowed to take leaves of $s$, not only internal vertices. However, if $t\in \ST$, there is a clear bijection between $\operatorname{Adm}(t)$ and the set of general admissible cuts of $\sk(t)$.
\end{Remark}

Let $t$ be a Schr\"oder tree with $\ell$ internal vertices whose root is denoted by $r$. For each $1\leq k\leq \ell$, define $\AdmT_k(t)$ to be the set of sequences $(c_1,\dots,c_{k-1})$, where, for each $j=1,\dots,k-1$, $c_j$ is a~general admissible cut of
\[
	\cR_{j-1}(\sk(t)) := R_{c_{j-1}}(R_{c_{j-2}}(\cdots R_{c_0}(\sk(t))\cdots)),
\]
with $\cR_0(\sk(t))=R_{c_0}(\sk(t)) := \sk(t)$ and $\cR_{k-1}(\sk(t)) =\{r\}$, such that any of
\[
	\ccP_{j}(\sk(t)):=P_{c_j}(R_{c_{j-1}}(R_{c_{j-2}}(\cdots R_{c_0}(\sk(t))\cdots)))
\]
is a nonempty forest of single-vertex trees.

\begin{Remark}Let $t\in \ST$ with root $r$ and $\tilde{c}= (c_1,\dots,c_{k-1})\!\in\! \AdmT_k(t)$. Since \mbox{$\cR_{k-1}(\sk(t)) = \{r\}$}, we have that $\tilde{c}$ is an ordered partition of the vertex set of $\sk(t)$ minus the root of $t$. Furthermore, since any $\ccP_j(\sk(t))$ is a nonempty forest of single-vertex trees, it is easy to see that $c_j$ is a subset of the set of leaves of $\cR_{j-1}(\sk(t))$ and $ \ccP_j(\sk(t)) = c_j,$ for any $1\leq j <k$.
\end{Remark}

The following lemma proves that the cardinality of the set $\AdmT_k(t)$ is given precisely by~$\omega_k(t)$.

\begin{Lemma}\label{lemma:bijection}Let $t$ be a Schr\"oder tree and $r$ be the roof of $t$. Then there is a bijection between $\AdmT_k(t)$ and the set of strictly order-preserving surjective functions $f\colon \sk(t)\to [k]$.
\end{Lemma}
\begin{proof}For any $(c_1,\dots,c_{k-1})\in \AdmT_k(t)$, we define the function $f\colon \sk(t)\to[k]$ by $f(r) = 1$, and $f(v) = k - j+1$ if $v\in c_{j}$.

We will show that the previous map is well defined. Indeed, observe that is $(c_1,\dots,c_k,\{r\})$ is an ordered partition of $\sk(t)$. Moreover, since every $c_j$ is a nonempty set, we have that $f$ is surjective. On the other hand, take $v$, $w$ vertices of $\sk(t)$ such that $v<w$, i.e., there is a~downward directed path from $v$ to $w$ in $\sk(t)$. The elements of the subsequence $(c_1,\dots,c_j)$ are consecutive prunings of $\sk(t)$ and these are taken from the leaves of $\cR_{j-1}(\sk(t))$ upward the root as $j$ increases. Then we have that $v\in c_i$, $w\in c_j$ and $j<i$, and in particular $f(v) = k-i+1 < k-j+1 = f(w)$. Hence $f$ is strictly order-preserving.

 It is clear that two different elements in $\AdmT_k(t)$ will produce two different functions. It remains to show that our correspondence is onto. To this end, for any $f\colon \sk(t)\to[k]$ strictly order preserving surjective function, define the sequence of sets $(c_1,\dots,c_{k-1})$ by
 \[
 c_j = \{ v\in \sk(t)\colon f(v) = k-j+1\},
 \]
 for any $1\leq j<k$. Now, observe that $c_1$ is a general admissible cut of $\sk(t)$ since $c_1$ is a subset of the set of leaves of $\sk(t)$. Actually, if $c_1$ contains a vertex $v$ that is not a leaf, there exists a~descendant of $v$, namely $w$ such that $k = f(v) < f(w)$ and this is a~contradiction. Hence $c_1$ is a~general admissible cut such that $\ccP_1(\sk(t)) = P_{c_1}(\sk(t)) = c_1$, i.e., $\ccP_{1}(\sk(t))$ is a~nonempty forest of single-vertex trees. Observe that the same argument can be applied to the tree $\cR_1(\sk(t))$ in order to show that $c_2$ is a general admissible cut such that $\ccP_{2}(\sk(t)) = c_2$ is a nonempty forest of single-vertex trees, and in general, $c_j$ is a general admissible cut such that $\ccP_j(\sk(t))= c_j$ is a nonempty forest of single-vertex trees, for any $j=1,\dots,k-1$. Since $f$ is order-preserving, we have that $f^{-1}(1) = \{r\}$. In addition since~$f$ is surjective, we have that $(c_1,\dots,c_{k-1},\{r\})$ is an ordered partition of the set of vertices of $\sk(t)$ with $\cR_{k-1}(\sk(t)) = \{r\}$. The above reasoning shows that $(c_1,\dots,c_{k-1})\in \AdmT_k(t)$ is a tuple such that their corresponding strictly order preserving surjective function is~$f$. Hence the map is onto, and the proof is complete.
\end{proof}

The main lemma of this section, which establishes the connection between the $\omega$ coefficients and the moments of random variables, is the following.

\begin{Lemma}
\label{lemma:main}
Let $(\A,\varphi)$ be a noncommutative probability space. Also, let $t \in \ST(n)$ and a~word $w=a_1\cdots a_n\in T_+(\A)$, for $n\geq1$. If $\tilde{\Phi}$ is the extension of $\varphi$ to a character as defined in~\eqref{eq:tildephi} and $\tilde\rho = \log^*(\tilde{\Phi})$, then we have
\begin{equation} \label{eq:tilderho10}
 \tilde\rho(t\otimes w) = \omega(\sk(t)) \varphi_{\pi(t)}(a_1,\dots,a_n).
\end{equation}
\end{Lemma}

\begin{proof}
Recall that
\[
	\tilde\rho(t\otimes w) = \sum_{k\geq 1} \frac{(-1)^{k+1}}{k} \hat\Phi^{*k} (t\otimes a_1\cdots a_n),
\]
where $\hat\Phi^{*k}$ is defined in \eqref{eq:iterated}. By \eqref{eq:muruacoeff}, it is enough to show that
\begin{equation}\label{eq:enough}
	\hat\Phi^{*k} (t\otimes a_1\cdots a_n) = \omega_k(\sk(t)) \varphi_{\pi(t)}(a_1,\dots,a_n),\qquad\mbox{for any }k\geq1.
\end{equation}
Let $k\geq 1$. First, we observe that according to the definition of $\delta_\A$, the $k$-fold iterated coproduct is of the form
\begin{equation}\label{eq:iteratedcop}
	\delta_\A^{[k]}(t\otimes a_1\cdots a_n) = \sum_{\mbox{\tiny{iterated admissible cuts of $t$}}}
	f_1\otimes\cdots \otimes f_k,
\end{equation}
where $f_1$ is a decorated subtree of $t$, each $f_2,\dots,f_k$ is a forest of decorated subtrees of $t$, and the sum is over iterated admissible cuts, i.e., doing admissible cuts on each iteration of the coproduct. Now observe that for any decorated tree, its evaluation on $\hat\Phi = \tilde\Phi - \varepsilon$ is equal to zero if the tree is empty or is not a~corolla. In particular, a~term $f_1\otimes\cdots \otimes f_k$ in the sum~\eqref{eq:iteratedcop} will produce a~zero contribution on $\tilde\Phi^{*k}$ if any of $f_1,\dots,f_k$ is the empty forest or is a~forest which contains a~tree that is not a corolla. Hence, the only terms that can produce a~non-zero contribution are such that $f_1$ is a corolla, and each of $f_2,\dots,f_k$ is a~forest of corollas. On the other hand, since corollas are associated to single-vertex trees via the skeleton map, by Remark~\ref{rem:generalcuts} the sum in~\eqref{eq:iteratedcop} is done precisely over the set $\AdmT_k(t)$.

Finally, by the definition of $\tilde\Phi$ in~\eqref{eq:tildephi}, the fact that $\hat\Phi$ is multiplicative, Definition \ref{def:pit}, and Lemma \ref{lemma:bijection}, the $k$-fold convolution product is equal to
\begin{align*}
	\hat\Phi^{*k}(t\otimes a_1\cdots a_n)
	&= \sum_{\mbox{\tiny{iterated admissible cuts of $t$}}} \hat\Phi(f_1)\cdots \hat\Phi(f_k)\\
	&= \sum_{\tilde{c}\in \AdmT_k(t)} \varphi_{\pi(t)}(a_1,\dots,a_n)= \omega_k(\sk(t)) \varphi_{\pi(t)}(a_1,\dots,a_n),
\end{align*}
where we get \eqref{eq:enough}.
\end{proof}

Now we are ready to prove Theorem~\ref{thm:mainthm1}.

\begin{proof}[Proof of Theorem \ref{thm:mainthm1}]
Let $\rho\colon T(T_+(\A))\to\mathbb{C}$ be the infinitesimal character associated to the monotone cumulants given by Theorem~\ref{thm:linkNCP}. By~\eqref{eq:rhoiota} and Lemma~\ref{lemma:main} we obtain
\begin{align*}
h_n(a_1,\dots,a_n) &= \rho(a_1\cdots a_n)
= \sum_{t\in \ST(n)} \tilde\rho(t\otimes a_1\cdots a_n)
\\ &= \sum_{t\in \ST(n)} \omega(\sk(t))\varphi_{\pi(t)} (a_1,\dots,a_n).\tag*{\qed}
\end{align*}\renewcommand{\qed}{}
\end{proof}
